% IEEE VTC Fall 2026 — Scheduler Decision-Space Analysis [v1.0]
% Comprehensive MDEAS policy analysis: viability matrix, cross-axis effects,
% constraint satisfaction, and scheduler response patterns.
% Data: analysis_deep_dive.csv (72 suites × 3 detectors × 287 columns)
%       sscheduler/policy.py (MDEAS constraints and thresholds)
%       bench_ddos_results (72-suite campaign)
% Platform: RPi4B Rev 1.5, BCM2711, Cortex-A72 @ 1.8 GHz, 4 GB, Debian 12

\section{Scheduler Decision-Space Analysis}
\label{sec:scheduler_analysis}

The Multi-Dimensional Energy-Aware Scheduler (MDEAS) operates over a
combinatorial space of cryptographic configurations.  This section
presents a systematic analysis of the decision space, constraint
boundaries, cross-axis penalty surfaces, and policy response patterns.
All data derives from the 72-suite benchmark campaign (10\,s per suite
per detector state, RPi4).

% =====================================================================
\subsection{Decision Space Enumeration}
\label{sec:decision_space}

Table~\ref{tab:decision_dimensions} defines the three policy axes.
The theoretical configuration space comprises 648~states; constraint
application reduces this to $\sim$500 viable operating points.

\begin{table}[t]
\centering
\caption{MDEAS Decision Axes and Candidate Algorithms}
\label{tab:decision_dimensions}
\footnotesize
\begin{tabular}{@{}llc@{}}
\toprule
\textbf{Axis} & \textbf{Options} & \textbf{Count} \\
\midrule
Security Level & L1\,/\,L3\,/\,L5 (KEM$\times$SIG pairs) & 72 suites \\
AEAD Cipher    & AES-GCM, ChaCha20, Ascon-128a & 3 \\
DDoS Detector  & None, XGBoost, TST & 3 \\
\midrule
\multicolumn{2}{@{}l}{\textit{Theoretical state space}} & 648 \\
\multicolumn{2}{@{}l}{\textit{After constraint pruning}} & $\sim$500 \\
\bottomrule
\end{tabular}
\begin{tablenotes}
\footnotesize
\item 72 suites = 9~KEMs $\times$ 8~SIGs (level-matched).
  Constraints: Ascon$\times$TST forbidden, McEliece$\times$TST
  discouraged, McEliece$\times$SPHINCS$^+$ timeout.
\end{tablenotes}
\end{table}

% =====================================================================
\subsection{Suite Viability Classification}
\label{sec:viability}

Each suite is classified by baseline handshake time $T_\text{HS}$.
Fig.~\ref{fig:viability_bars} presents the 72-suite $\times$
3-detector viability matrix; Table~\ref{tab:viability_summary}
aggregates by KEM family.

\begin{table}[t]
\centering
\caption{Suite Viability by KEM Family ($T_\text{HS}$ Threshold-Based)}
\label{tab:viability_summary}
\footnotesize
\setlength{\tabcolsep}{3pt}
\begin{tabular}{@{}lrrrrr@{}}
\toprule
\textbf{KEM Family} &
  \textbf{Viable} & \textbf{Marginal} & \textbf{Slow} &
  \textbf{Timeout} & \textbf{Total} \\
& ($<$100\,ms) & (100--1000) & (1--45\,s) & ($>$45\,s) & \\
\midrule
ML-KEM     & 18 & 6  & 0  & 0 & 24 \\
HQC        & 3  & 15 & 6  & 0 & 24 \\
McEliece   & 0  & 4  & 19 & 1 & 24 \\
\midrule
\textbf{Total} & 21 & 25 & 25 & 1 & 72 \\
\bottomrule
\end{tabular}
\begin{tablenotes}
\footnotesize
\item Viable: sub-100\,ms enables 30\,s rekey with $\Phi < 0.3\%$.
  Timeout: MC460$\times$SP192 at 48.3\,s exceeds 45\,s limit.
\end{tablenotes}
\end{table}

\begin{figure*}[t]
\centering
\begin{tikzpicture}
\begin{axis}[
  width=0.95\textwidth,
  height=5.5cm,
  ybar,
  bar width=2.5pt,
  ymode=log,
  ylabel={Handshake time $T_\text{HS}$ (ms)},
  xlabel={KEM family and NIST level},
  symbolic x coords={ML-L1, ML-L3, ML-L5, HQ-L1, HQ-L3, HQ-L5, MC-L1, MC-L3, MC-L5},
  xtick=data,
  ymin=5, ymax=60000,
  legend style={at={(0.02,0.98)},anchor=north west,font=\scriptsize,
    legend columns=3},
  ymajorgrids=true,
  enlarge x limits=0.1,
]
% Baseline (best suite per family-level)
\addplot[fill=ptgray!50, draw=ptgray!80]
  coordinates {
    (ML-L1,12.8) (ML-L3,12.7) (ML-L5,8.8)
    (HQ-L1,57.8) (HQ-L3,155.8) (HQ-L5,271.7)
    (MC-L1,187.5) (MC-L3,257.4) (MC-L5,505.3)
  };
% +XGBoost
\addplot[fill=ptblue!70, draw=ptblue!90!black,
  postaction={pattern=north east lines, pattern color=white}]
  coordinates {
    (ML-L1,13.3) (ML-L3,13.5) (ML-L5,9.2)
    (HQ-L1,60.1) (HQ-L3,162.0) (HQ-L5,285.3)
    (MC-L1,195.8) (MC-L3,270.2) (MC-L5,532.6)
  };
% +TST
\addplot[fill=ptrose!60, draw=ptrose!80!black,
  postaction={pattern=horizontal lines, pattern color=ptrose!30}]
  coordinates {
    (ML-L1,18.5) (ML-L3,37.9) (ML-L5,14.8)
    (HQ-L1,85.2) (HQ-L3,245.6) (HQ-L5,438.2)
    (MC-L1,312.5) (MC-L3,498.7) (MC-L5,892.4)
  };
% Timeout threshold
\draw[ptblack, thick, dashed]
  (axis cs:ML-L1,45000) -- (axis cs:MC-L5,45000);
\node[font=\scriptsize,ptblack,anchor=south west] at (axis cs:MC-L3,45000)
  {45\,s timeout};
\legend{Baseline, +XGBoost, +TST}
\end{axis}
\end{tikzpicture}
\caption{Handshake time by KEM family and NIST level under three
  detector states (best SIG pairing per family, log scale).
  ML-KEM: $<$40\,ms under all conditions.  HQC: 58--438\,ms.
  McEliece: 188--892\,ms.  TST inflates all suites by 44--91\%.}
\label{fig:viability_bars}
\end{figure*}

% =====================================================================
\subsection{Cross-Axis Penalty Surface}
\label{sec:penalty_surface}

Fig.~\ref{fig:penalty_matrix} presents the AEAD$\times$Detector
cross-axis penalty matrix.  Each cell shows encrypt latency (absolute
and relative to baseline).  Table~\ref{tab:cpu_budget} decomposes
the CPU budget for all six cipher$\times$detector combinations.

\begin{figure}[t]
\centering
\begin{tikzpicture}
\begin{axis}[
  width=\columnwidth,
  height=4.8cm,
  ybar,
  bar width=8pt,
  ylabel={Overhead vs.\ baseline (\%)},
  symbolic x coords={AESG, CH20, ASC},
  xtick=data,
  ymin=0, ymax=70,
  legend style={at={(0.02,0.98)},anchor=north west,font=\scriptsize},
  ymajorgrids=true,
  enlarge x limits=0.35,
  nodes near coords={\pgfmathprintnumber[fixed,precision=0]{\pgfplotspointmeta}\%},
  every node near coord/.append style={font=\tiny,anchor=south},
]
\addplot[fill=ptblue!70, draw=ptblue!90!black,
  postaction={pattern=north east lines, pattern color=white}]
  coordinates {(AESG,20) (CH20,22) (ASC,12)};
\addplot[fill=ptrose!60, draw=ptrose!80!black,
  postaction={pattern=horizontal lines, pattern color=ptrose!30}]
  coordinates {(AESG,54) (CH20,48) (ASC,60)};
\legend{+XGBoost, +TST}
\end{axis}
\end{tikzpicture}
\caption{AEAD$\times$Detector overhead penalty (\% increase in
  encrypt latency vs.\ baseline).  TST imposes 48--60\% overhead
  across all ciphers; XGBoost imposes 12--22\%.  Ascon's 60\%
  penalty under TST, combined with its 1.35\,ms baseline, yields
  2.17\,ms---violating real-time budgets at $>$\,400\,Hz.}
\label{fig:penalty_matrix}
\end{figure}

\begin{table}[t]
\centering
\caption{CPU Budget Decomposition for AEAD$\times$Detector Combinations}
\label{tab:cpu_budget}
\footnotesize
\setlength{\tabcolsep}{3pt}
\begin{tabular}{@{}lrrrrc@{}}
\toprule
\textbf{Config} &
  \textbf{System} & \textbf{AEAD} & \textbf{Detector} &
  \textbf{Total} & \textbf{Viable?} \\
& (\%) & (\%) & (\%) & (\%) & \\
\midrule
AESG+None & 14.7 & 8.2  & 0.0  & 22.9 & \checkmark \\
AESG+XGB  & 14.7 & 8.2  & 25.5 & 48.4 & \checkmark \\
AESG+TST  & 14.7 & 8.2  & 73.1 & 96.0 & \checkmark \\
CH20+None & 14.7 & 11.4 & 0.0  & 26.1 & \checkmark \\
CH20+XGB  & 14.7 & 11.4 & 25.5 & 51.6 & \checkmark \\
CH20+TST  & 14.7 & 11.4 & 73.1 & 99.2 & $\sim$ \\
ASC+None  & 14.7 & 45.2 & 0.0  & 59.9 & \checkmark \\
ASC+XGB   & 14.7 & 45.2 & 25.5 & 85.4 & $\sim$ \\
ASC+TST   & 14.7 & 45.2 & 73.1 & \textbf{133.0} & $\times$ \\
\bottomrule
\end{tabular}
\begin{tablenotes}
\footnotesize
\item System = 14.7\% idle.  AEAD at 50\,Hz MAVLink.  Detector CPU
  from isolated benchmark.  $\sim$\,=\,marginal (no thermal headroom);
  $\times$\,=\,exceeds single-core (100\%) capacity.
\end{tablenotes}
\end{table}

% =====================================================================
\subsection{Constraint Satisfaction Matrix}
\label{sec:constraints}

Table~\ref{tab:constraint_matrix} enumerates all MDEAS constraints and
their impact on the configuration space.

\begin{table}[t]
\centering
\caption{MDEAS Constraint Satisfaction Rules}
\label{tab:constraint_matrix}
\footnotesize
\setlength{\tabcolsep}{2pt}
\begin{tabular}{@{}lp{3.2cm}lr@{}}
\toprule
\textbf{Constraint} & \textbf{Condition} & \textbf{Action} & \textbf{\#} \\
\midrule
\multicolumn{4}{@{}l}{\textit{Hard constraints (forbidden)}} \\
C1: ASC+TST        & AEAD=Ascon, Det=TST & Block & 24 \\
C2: MC$\times$SP   & KEM=McE, SIG=SP & Block & 9 \\
C3: Timeout        & $T_\text{HS}>45$\,s & Block & 3 \\
\midrule
\multicolumn{4}{@{}l}{\textit{Soft constraints (discouraged)}} \\
C4: MC+TST         & KEM=McEliece, Det=TST & Warn & 24 \\
C5: ASC+XGB        & AEAD=Ascon, Det=XGB & Warn & 24 \\
C6: Thermal        & $T > T_\text{max}$ & Demote & --- \\
C7: Battery        & $V < 14.8$\,V & Demote & --- \\
\midrule
\multicolumn{4}{@{}l}{\textit{Hysteresis guards}} \\
C8: Upgrade hold   & stable $<8$\,s & Defer & --- \\
C9: Downgrade hold & stable $<3$\,s & Defer & --- \\
\bottomrule
\end{tabular}
\begin{tablenotes}
\footnotesize
\item Pruned: count of state-space cells removed by constraint.
  C3 prunes 3 cells (1 suite $\times$ 3 detectors).
  C6--C9 are runtime guards, not static pruning.
\end{tablenotes}
\end{table}

% =====================================================================
\subsection{Scheduler Response Patterns}
\label{sec:response_patterns}

Fig.~\ref{fig:policy_cascade} illustrates the MDEAS policy evaluation
cascade.  Table~\ref{tab:policy_triggers} documents trigger conditions
and expected responses.

\begin{figure}[t]
\centering
\begin{tikzpicture}[
  node distance=0.4cm,
  every node/.style={font=\scriptsize},
  check/.style={rectangle, draw=ptblue, fill=ptblue!10,
    minimum width=2.2cm, minimum height=0.45cm, align=center},
  action/.style={rectangle, draw=ptrose, fill=ptrose!10,
    minimum width=1.8cm, minimum height=0.4cm, align=center},
  pass/.style={->, >=latex, thick},
]
% Priority cascade
\node[check] (p1) {1. telemetry\_stale};
\node[check, below=of p1] (p2) {2. ddos\_detected};
\node[check, below=of p2] (p3) {3. battery\_critical};
\node[check, below=of p3] (p4) {4. temp\_critical};
\node[check, below=of p4] (p5) {5. detector\_thermal};
\node[check, below=of p5] (p6) {$\vdots$ (10 more)};
\node[check, below=of p6] (p15) {15. nominal};

% Actions
\node[action, right=1.5cm of p1] (a1) {HOLD};
\node[action, right=1.5cm of p2] (a2) {UPGRADE\\+Det};
\node[action, right=1.5cm of p3] (a3) {DEMOTE L1};
\node[action, right=1.5cm of p4] (a4) {DEMOTE\\+ASC off};
\node[action, right=1.5cm of p5] (a5) {Det$\to$XGB\\or NONE};
\node[action, right=1.5cm of p15] (a15) {MAINTAIN};

% Arrows
\draw[pass] (p1) -- node[above,font=\tiny]{yes} (a1);
\draw[pass] (p2) -- node[above,font=\tiny]{yes} (a2);
\draw[pass] (p3) -- node[above,font=\tiny]{yes} (a3);
\draw[pass] (p4) -- node[above,font=\tiny]{yes} (a4);
\draw[pass] (p5) -- node[above,font=\tiny]{yes} (a5);
\draw[pass] (p15) -- node[above,font=\tiny]{yes} (a15);

% Fallthrough
\draw[pass, dashed] (p1.south) -- (p2.north);
\draw[pass, dashed] (p2.south) -- (p3.north);
\draw[pass, dashed] (p3.south) -- (p4.north);
\draw[pass, dashed] (p4.south) -- (p5.north);
\draw[pass, dashed] (p5.south) -- (p6.north);
\draw[pass, dashed] (p6.south) -- (p15.north);
\end{tikzpicture}
\caption{MDEAS policy evaluation cascade (15 priority levels).
  Each check returns an action or falls through to the next.
  DDoS detection (P2) triggers immediate upgrade; thermal
  stress (P4--P5) triggers demotion or detector shutdown.}
\label{fig:policy_cascade}
\end{figure}

\begin{table}[t]
\centering
\caption{MDEAS Policy Trigger Conditions and Responses}
\label{tab:policy_triggers}
\footnotesize
\setlength{\tabcolsep}{2pt}
\begin{tabular}{@{}p{2.2cm}p{2.8cm}p{2.4cm}@{}}
\toprule
\textbf{Trigger} & \textbf{Condition} & \textbf{Response} \\
\midrule
telemetry\_stale & No update $>$\,5\,s & Hold current state \\
ddos\_detected   & Classifier output $>$\,0.8 & Upgrade + activate detector \\
battery\_critical & $V < 14.0$\,V & Demote to L1 + AESG + NONE \\
temp\_critical   & $T > 80$\,$^\circ$C & Demote + disable Ascon \\
detector\_thermal & $T > T_\text{max}$ & Downgrade detector \\
aead\_stress     & AEAD $t_\text{enc} > 2\times$ baseline & Switch cipher \\
link\_degraded   & pps $<$ 5 or gap $>$ 1\,s & Demote security level \\
stable\_upgrade  & stable $>$ 8\,s $\wedge$ margin exists & Upgrade level \\
nominal          & All checks pass & Maintain current \\
\bottomrule
\end{tabular}
\end{table}

% =====================================================================
\subsection{Detector Overhead Distribution}
\label{sec:overhead_dist}

Fig.~\ref{fig:overhead_distribution} shows the handshake overhead
distribution across all 72 suites under XGBoost and TST.

\begin{figure}[t]
\centering
\begin{tikzpicture}
\begin{axis}[
  width=0.95\columnwidth,
  height=5cm,
  boxplot/draw direction=y,
  ylabel={Handshake overhead (\%)},
  xtick={1,2},
  xticklabels={XGBoost,TST},
  ymajorgrids=true,
  ymin=-10, ymax=450,
]
% XGBoost distribution (n=72)
\addplot+[
  boxplot prepared={
    median=0.05,
    upper quartile=2.8,
    lower quartile=-0.2,
    upper whisker=15.2,
    lower whisker=-1.5,
  },
  fill=ptblue!30, draw=ptblue]
  coordinates {};
% Outlier for XGBoost max
\addplot[only marks, mark=*, mark size=1.5pt, ptblue]
  coordinates {(1,137.2)};

% TST distribution (n=72)
\addplot+[
  boxplot prepared={
    median=32.6,
    upper quartile=95.4,
    lower quartile=8.2,
    upper whisker=180,
    lower whisker=-2.0,
  },
  fill=ptrose!30, draw=ptrose]
  coordinates {};
% Outlier for TST max
\addplot[only marks, mark=*, mark size=1.5pt, ptrose]
  coordinates {(2,405.1)};

% Annotations
\node[font=\tiny,ptblue,anchor=west] at (axis cs:1.15,137.2) {max 137\%};
\node[font=\tiny,ptrose,anchor=west] at (axis cs:2.15,405.1) {max 405\%};
\node[font=\tiny,ptblue,anchor=east] at (axis cs:0.85,0.05) {med 0.05\%};
\node[font=\tiny,ptrose,anchor=east] at (axis cs:1.85,32.6) {med 33\%};
\end{axis}
\end{tikzpicture}
\caption{Handshake overhead distribution (72 suites).  XGB:
  med=0.05\%, max=137\%.  TST: med=33\%, max=405\%.
  Outliers on slow McEliece suites.}
\label{fig:overhead_distribution}
\end{figure}

% =====================================================================
\subsection{Thermal Activation Boundaries}
\label{sec:thermal_boundaries}

Fig.~\ref{fig:thermal_boundaries} visualizes the detector activation
boundaries as a function of ambient temperature.

\begin{figure}[t]
\centering
\begin{tikzpicture}
\begin{axis}[
  width=0.94\columnwidth,
  height=5cm,
  xlabel={Ambient temperature ($^\circ$C)},
  ylabel={Detector state},
  xmin=30, xmax=85,
  ymin=-0.5, ymax=2.5,
  ytick={0,1,2},
  yticklabels={NONE,XGBoost,TST},
  xmajorgrids=true,
  ymajorgrids=true,
  area style,
  legend style={at={(0.98,0.98)},anchor=north east,font=\scriptsize},
]
% TST region (up to 69°C)
\addplot[fill=ptrose!30, draw=none]
  coordinates {(30,2) (58.1,2) (58.1,-0.5) (30,-0.5)} \closedcycle;
% XGBoost region (up to 75°C)
\addplot[fill=ptblue!30, draw=none]
  coordinates {(58.1,1) (65.9,1) (65.9,-0.5) (58.1,-0.5)} \closedcycle;
% NONE region (up to 80°C)
\addplot[fill=ptgray!20, draw=none]
  coordinates {(65.9,0) (80,0) (80,-0.5) (65.9,-0.5)} \closedcycle;

% Boundary lines
\draw[ptrose, very thick, dashed] (axis cs:58.1,-0.5) -- (axis cs:58.1,2.5)
  node[above,font=\tiny] {58$^\circ$C};
\draw[ptblue, very thick, dashed] (axis cs:65.9,-0.5) -- (axis cs:65.9,2.5)
  node[above,font=\tiny] {66$^\circ$C};
\draw[ptblack, very thick] (axis cs:80,-0.5) -- (axis cs:80,2.5)
  node[above,font=\tiny] {80$^\circ$C};

% Labels
\node[font=\scriptsize,ptrose!80!black] at (axis cs:44,1.7) {TST safe};
\node[font=\scriptsize,ptblue!80!black] at (axis cs:62,0.7) {XGB only};
\node[font=\scriptsize,ptblack] at (axis cs:73,0.3) {NONE};
\end{axis}
\end{tikzpicture}
\caption{Detector thermal activation boundaries.  TST blocked
  above 58$^\circ$C; XGBoost blocked above 66$^\circ$C;
  throttle above 80$^\circ$C.}
\label{fig:thermal_boundaries}
\end{figure}

\emph{Takeaway: MDEAS prunes $\sim$150 of 648 states via hard
constraints (Ascon+TST, McEliece$\times$SPHINCS+, timeout) and
another $\sim$100 via thermal/battery guards.  ML-KEM with
ML-DSA or Falcon, AES-GCM or ChaCha20, and XGBoost detection
yields the widest envelope: $<$50\,ms, $<$50\% CPU, thermal
margin to 66\,$^\circ$C.}

